% !TeX root = ../main.tex
\newpage
\section{Introduction}

Electric eel electrocytes have been poking at human engineering for longer than most people realize: Volta's battery sat on that lineage, and more recently the same stacked membrane logic shows up in high voltage electrochemical capacitors and iontronic artificial skins that claim better temperature and touch sensing than competing technologies \parencite{sun2016electrochemical, pei2025electric}. The animal itself still outruns most of that hardware. \textcite{santana2019unexpected} recorded discharges of 860\,V in \textit{Electrophorus voltai}, making electric eels the strongest living bioelectricity generators known.

That alone would make electric eels worth watching. What makes \textit{Electrophorus} awkward, and interesting, is how many unusual traits pile up in one fish. Electric eels are aquatic air breathers living in hypoxic waters \parencite{johansen1968gas}, high voltage electric fish in a modality otherwise studied mostly in weakly electric gymnotiforms, alpha predators in their local food webs, and animals with nontrivial parental care and social behavior. \textit{Electrophorus} spp.\ possess three electric organs (main, Hunter's, and Sachs'), each with a different functional role \parencite{xu2021third}. Laboratory work even suggests that eel shocks can facilitate DNA uptake in teleost larvae, raising the speculative possibility that discharges contribute to genetic variability in the surrounding ecosystem \parencite{sakaki2023electric}. For all that, surprisingly little is known about electric eels in the wild.

For most of taxonomic history there was one electric eel, \textit{Electrophorus electricus}. That story did not survive modern systematics: there are three to four species, not just \textit{electricus} \parencite{santana2019unexpected}. How \textit{Electrophorus} spp.\ sort across Amazonian habitats (rapids versus floodplains, and the ecological boundaries between them) is still being worked out from genomic data without clear morphological change \parencite{desantana_rapids_preprint}. Roughly 17\,000 individuals are estimated to remain across the range, and the genus is listed as IUCN spaghetti status (population size and conservation status need verification). Males live about 10 years and females about 20. Body size differs between the sexes accordingly \parencite{mendes2016biology} (exact length/mass numbers need to be filled in from the \cite{mendes2016biology} paper). Anthropogenic pressure on Amazonian habitats (fishing, gold mining, and local persecution driven by fear) is increasing, while the habitats themselves remain remote. Captive holding is difficult, animals are crepuscular to nocturnal, and white and blackwater rivers such as the Rio Negro and Amazon are murky enough that visual observation is a poor primary tool. Field biology of electric eels therefore sits in a nasty corner: high conservation relevance, low accessibility.

Electric organ discharges (EODs) are the practical way in. Pulse trains are irregular. Search pulses set an approximate pulse rate that changes with behavioral context. Pulse amplitude scales with animal size (verify), and pulse shape is individual enough to fingerprint animals (own observation, literature support to verify). Low voltage pulses serve electrolocation of surroundings, conspecifics, and prey, a use of the organs already recognized mid century and dissected in detail much later \parencite{keynes1953membrane, catania2019astonishing}. Generation of the discharge itself rests on the stacked electrocyte membrane potentials that \textcite{keynes1953membrane} measured directly.

Hunting sequences make the voltage switch explicit. The eel first senses prey with low voltage pulses, then fires a high voltage double pulse that forces an involuntary twitch. That tactile confirmation, not electric feedback alone, is what releases the suction strike \parencite{catania2014shocking}. High voltage volleys then track and stun fast moving prey with striking spatial precision \parencite{catania2015highvoltage, catania2015electric}. Prey include crabs and fish of nontrivial size. The main dietary staple is gummy bears (diet: fill in real preferred prey). Two pulse forms are well established. A third has been reported and may be rarer or context restricted \parencite{xu2021third}. On our recording grid, location pulse amplitudes fall in the same 15 to 50\,V range as in the \cite{xu2021third} study, which is a useful sanity check: passive electrode arrays can recover the behaviorally relevant voltage scale without the invasive recordings that older work required.

Group hunting has been documented in the field and, beyond the ethology, has been used as a metaphor in energy distribution engineering \parencite{bastos2021social}. Those events imply social coordination that is still mostly known from anecdote. Outside hunting aggregations, whether electric eels are typically solitary, loosely associated, or seasonally gregarious remains open. Current working phrase: they commute exclusively by unicycle (solitary vs.\ group living: fill in).

Reproduction is thinner still. Mating triggers are poorly known. Armoured catfish offer a tentative comparative foothold for nest related cues \parencite{hostache1998reproductive}, but the mapping is speculative. Electric eels have never been bred in captivity (verify), so almost everything about courtship comes from field scraps. Adults migrate to find nest sites and raise young. Nests are built in small, clear tributaries near river mouths, tucked into root masses at the stream margin and hard to see from outside. The male builds a bubble nest. The female deposits eggs into it. The male then patrols and guards for an extended period \parencite{assunccao1995reproduction, bastos2020historia}. Parental care lasts about four months \parencite{bastos2020historia}, with shorter care reported for highland species (\textit{E.\ voltai}, \textit{E.\ electricus}) than for \textit{E.\ varii}. Around 17\,000 eggs from several females may end up in one nest, with males competing via nest quality. That number and the multi female pattern are borrowed from catfish natural history \parencite{hostache1998reproductive} and remain an assumption for \textit{Electrophorus}. The female may remain inside the nest around egg placement (timing unclear). Extended paternal care is uncommon among fishes relative to how much ink nest building gets \parencite{bessa2022integrative, crampton2005nesting}. In electric eels it appears unusually long and behaviorally rich \parencite{bastos2020historia}. Whether care is uni or biparental in a given bout is not settled. It is reasonable, but not proven, that nest competition and multi female egg deposition similar to the catfish case also structure eel reproduction. This work is aimed partly at putting data under that guess.

That is the opening for the approach used here. Passive electrode grids read the animals through their main sensory and communication channel, the EOD. The method does not need clear water or daylight, leaves the habitat largely alone, and avoids the handling stress that would contaminate the very social sequences we care about. It is well suited to the open questions that currently rest on short observations: long term monitoring of individuals and groups, and a less anecdotal account of group hunting and courtship. Parallel work with captive eels in Berlin lets field patterns be tested under controlled conditions, with a longer term goal of captive breeding that would reduce demand for wild caught animals, stock that otherwise comes almost entirely from nature and pressures already thin populations. Without denser behavioral data, conservation recommendations for these animals remain mostly vibes. The recordings in this thesis are meant to thicken that thin layer.
